\chapter{实证分析与结果对比}

第五章基于经典金融文献中的基准参数案例，系统完成了 3/2 随机波动率模型下九类欧式期权定价方法的精度、效率与数值稳定性评价，明确了不同方法在理想化参数环境下的理论性能边界与适用场景。但需要指出的是，现有学术研究中的基准测试案例多为平稳市场环境下的理想化参数设定，参数取值无极端波动、无突发市场冲击，虽能完成方法的理论有效性验证，却无法反映真实金融市场、尤其是尾部风险集中的高波动市场中，模型参数的时变特征与定价方法的实际工程适配性。

加密货币市场作为全球波动水平最高的大类资产市场之一，具有标的价格波动剧烈、杠杆交易占比高、极端行情下尾部风险快速释放的核心特征，其期权市场的波动率微笑形态在极端事件中会出现快速且剧烈的形变，为检验随机波动率模型的市场解释力、参数动态特征与定价方法的极端场景鲁棒性提供了天然的实证实验场景。本章以 2025 年 10 月 10 日至 11 日的 BTC 市场闪崩事件为研究对象，将 3/2 模型的定价分析从学术基准算例拓展至真实极端市场环境，完成模型的全周期参数校准与定价方法的实证再检验。

本章的核心研究目标分为三个层面：第一，识别 3/2 模型核心参数在崩盘事件前、崩盘期间、市场恢复期三个阶段的时变规律，明确极端行情对模型参数的冲击特征，验证 3/2 模型对加密货币期权市场波动率微笑的拟合能力；第二，基于真实市场校准得到的参数构建全新的测试案例，重复第五章的定价方法性能对比，检验不同方法在极端市场场景下的数值稳定性、定价精度与计算效率，弥补学术基准案例与真实市场环境的脱节问题；第三，明确 3/2 模型在加密货币期权定价中的实践适配性，为极端行情下的衍生品定价、风险管理提供可落地的实证依据与方法选择指引。

本章的研究逻辑严格遵循 “长期固定参数估计 - 横截面时变参数校准 - 定价方法极端场景再检验” 的完整链条，所有实证分析均基于 Deribit 交易所的真实逐笔成交数据完成。本章后续内容安排如下：6.2 节明确研究事件的时间窗口划分、数据来源与预处理规则；6.3 节构建 3/2 模型的两步校准框架，说明长期参数估计与横截面参数校准的理论基础与实现逻辑；6.4 节完成崩盘事件全周期期权成交数据的描述性统计分析，明确校准样本的筛选依据；6.5 节呈现全周期的参数校准结果，分析极端行情下模型核心参数的时变特征与模型拟合效果；6.6 节基于校准得到的真实市场参数，完成定价方法的性能再检验，并与第五章的学术案例结果进行对比分析；6.7 节对本章的核心实证结论进行总结。

\section{数据准备与处理}

为了验证 3/2 模型在极端市场情境下的鲁棒性，本研究构建了两类数据集：定价测试数据：选取文献中公认的极端场景参数作为基准，并结合 Euler/Milstein 离散化方案生成的极端情景模拟数据，用于评估 $Sv32Fft$ 算法在平值点及临近到期时的数值精度。校准实证数据：波动率指标：获取自 2025 年 4 月 13 日至 2025 年 10 月 11 日期间 Deribit 交易所的 DVOL 指数 K 线数据，用于捕捉市场全局波动率水平。期权逐笔数据：重点选取 2025 年 10 月 10 日 Bitcoin 发生剧烈波动前后的三个关键时间段：冲击前阶段（20:55:00 - 21:00:00 UTC）暴跌阶段（21:15:00 - 21:20:00 UTC）：此阶段 Bitcoin 价格从 113,000 区域迅速回落至 101,500 附近，展现了剧烈的下行冲击。修复阶段（23:05:00 - 23:10:00 UTC）

\section{校准方法与实现路径}

本研究采用基于傅里叶变换（FFT）的快速定价框架，利用 $Sv32Fft$ 类对 3/2 模型参数进行逆向校准。校准目标函数定义为模型隐含波动率与市场观测隐含波动率之间的均方误差（MSE）最小化：$$\min_{\Theta} \sum_{i=1}^N w_i \left( \sigma_{market, i} - \sigma_{model, i}(\Theta) \right)^2$$其中状态向量 $\Theta = \{\sigma, \kappa, \theta, \epsilon, \rho\}$。在实现过程中，利用前述章节推导的合流超几何函数 $_1F_1$ 及其对数 Gamma 变换，有效解决了校准过程中频繁调用特殊函数导致的计算溢出问题。针对加密市场的高频特性，校准过程配合了分布式参数寻优策略，以确保在短时间内完成对全期限波动率曲面的动态拟合。

\section{结果分析与讨论}

\subsection{波动率微笑的演变特征}
校准结果显示，在不同到期期限（如 6 天与 20 天）下，3/2 模型对市场“波动率微笑”的拟合表现出显著的时间敏感性。在 10 月 10 日的暴跌过程中（Crash 阶段），近月期权（6-Day Maturity）的 ATM 隐含波动率从 41.9\% 飙升至 55.4\%，且曲面的左侧偏斜（Left Skew）急剧增大，反映出市场对进一步下行风险的极度恐慌。

\subsection{3/2 模型与纯扩散模型的对比}
通过对比“纯扩散拟合（Pure Diffusion Fit）”与包含随机波动率项的校准结果可以发现：极端行情拟合：在冲击时刻，纯扩散模型无法解释深虚值（OTM）看跌期权的高溢价，导致拟合曲线在尾部出现显著背离。参数动态：校准得到的 $vov$（波动率之波动率）在 Crash 阶段显著升高，同时相关系数 $\rho$ 呈现出强烈的负相关特征（接近 -0.99），这证明了 3/2 模型在刻画“价格下跌触发波动率暴涨”这一负反馈机制上的卓越性能。恢复阶段特征：随着市场进入修复期，3/2 模型的参数表现出良好的回归特性，均值回归速度 $\kappa$ 的估计值反映了加密市场较之传统市场更快的风险出清节奏。

\subsection{鲁棒性总结}
实证分析表明，本研究所提出的 $Sv32Fft$ 框架在面对 Bitcoin 超过 10\% 的单日瞬时跌幅时，依然能够保持数值计算的稳定性，未出现特征函数震荡导致的定价失效。这不仅验证了 3/2 模型的理论优越性，也证明了针对修正式贝塞尔函数及其导数所做的数学优化在实务操作中的必要性。